\chapter[Hausandacht]{Hausandacht}
\quellseite{10}
Der berühmte Geschichtsforscher Leopold von Ranke pflegte jeden Morgen mit einem kurzen Worte aus der heiligen Schrift sich für die Arbeit des Tages zu stärken. Sein Sohn, ein Prediger in Preußen, theilte vor einiger Zeit auf einer Synode Folgendes darüber mit: „Viele Jahre haben mein Bruder und ich, ehe wir ins Gymnasium gingen, unserm Vater aus einem kleinen Büchlein, die ‚Dreifältige Schnur‘ genannt, drei Sprüche der heiligen Schrift vorgelesen. Sonntags erweiterte sich diese Morgenandacht, indem wir dem Vater das Evangelium und die Epistel des Sonntags vorlasen. Auch die Feste, die kirchlichen sowohl wie die häuslichen, wurden durch das Lesen eines Psalmes oder einer andern Bibelstelle geweiht. Alle Jahre einmal wurde das heilige Abendmahl im Hause gefeiert. Da steht mir eine Feier besonders vor Augen: Oberhofprediger Dr. Strauß, der mich confirmirt hatte, hatte das Abendmahl ausgetheilt. Wir hatten auf den Knien das Dankgebet gesprochen; nun erhob sich mein Vater und sagte, in einem evangelischen Hause wäre der Hausvater auch der Hauspriester, und nun rühmte er die Güte und Gnade seines Heilandes. Aus dieser Hausandacht ist uns ein unverlierbarer Segen erwachsen.“
